\begin{quote}
On a day nearly 4,000 years ago, in a town called Menet Khufu bordering the thin ribbon of the Nile, a master scribe sketched out the hieroglyphs that told the story of his lord's life—and in so doing he opened the recorded history of cryptology. His was not a system of secret writing as the modern world knows it; he used no fully developed code of hieroglyphic symbol substitutions. His inscription, carved about 1900 B.C. into the living rock in the main chamber of the tomb of the nobleman Khnumhotep II, merely uses some unusual hieroglyphic symbols here and there in place of the more ordinary ones. Most occur in the last 20 columns of the inscription's 222, in a section recording the monuments that Khnumhotep had erected in the service of the pharaoh Amenemhet II.

The intention was not to make it hard to read the text. It was to impart a dignity and authority to it, perhaps in the same way that a government proclamation will spell out ``In the year of Our Lord One thousand eight hundred and sixty three'' instead of just writing ``1863.'' The anonymous scribe may also have been demonstrating his knowledge for posterity. Thus the inscription was not secret writing, but it incorporated one of the essential elements of cryptography: a deliberate transformation of the writing. It is the oldest text known to do so.

The transformations occur in funerary formulas, in a hymn to Thoth, in a chapter of the Book of the Dead, on the sarcophagus of the pharaoh Seti I, in royal titles displayed in Luxor, on the architrave of the Temple of Luxor, on stele, in laudatory biographic inscriptions. There is nothing meant to be concealed in all this; indeed, many of the statements are repeated in ordinary form right next to the altered ones. Why, then, the transformations? Sometimes for essentially the same reason as in Khnumhotep's tomb: to impress the reader. Occasionally for a calligraphic or decorative effect; rarely, to indicate a contemporary pronunciation; perhaps even for a deliberate archaism as a reaction against foreign influence.

But many inscriptions are tinctured, for the first time, with the second essential for cryptology—secrecy. In a few cases, the secrecy was intended to increase the mystery and hence the arcane magical powers of certain religious texts. But the secrecy in many more cases resulted from the understandable desire of the Egyptians to have passersby read their epitaphs and so confer upon the departed the blessings written therein.

In Egypt, with its concentration upon the afterlife, the number of these inscriptions soon proliferated to such an extent that the attention and the goodwill of visitors flagged. To revive their interest, the scribes deliberately made the inscriptions a bit obscure. They introduced the cryptographic signs to catch the reader's eye, make him wonder, and tempt him into unriddling them—and so into reading the blessings. It was a sort of Madison Avenue technique in the Valley of the Kings. But the technique failed utterly. Instead of interesting the readers, it evidently destroyed even the slightest desire to read the epitaphs, for soon after the funerary cryptography was begun, it was abandoned.
\end{quote}
